\input{../common/header.tex}
\subtitle{Start programming in Python}
\date{2020-04-02}
\input{../common/title.tex}

\begin{frame}
	\tableofcontents
\end{frame}

\section{Organizational matters}
\begin{frame}[fragile]{Organizational matters}
	\begin{itemize}
		\item Please participate in writing today's lecture notes:
            \href{https://yourpart.eu/p/lecture-scientific-computing03-notes}{https://yourpart.eu/p/lecture-scientific-computing03-notes}\pause
        \item Glossary:\\
            \href{https://yourpart.eu/p/lecture-scientific-computing-glossary}{https://yourpart.eu/p/lecture-scientific-computing-glossary}
	\end{itemize}
\end{frame}

\section{Leftover from last week}

\begin{frame}[fragile]{Using Jupyter Notebooks}
	
	If you want to use conda environments in your notebooks, do the following:
	\begin{itemize}
		\item Install \textit{ipykernel} in your environment
		
		\begin{verbatim}
		conda activate <conda-environment>
		conda install ipykernel
		conda deactivate
		\end{verbatim}
		\item Additionally, install \textit{nb\_conda\_kernels} in \textit{base} if you run your jupyter notebook from base:
		\begin{verbatim}
		conda activate base    # could be also some other environment
		conda install nb_conda_kernels
		\end{verbatim}
	\end{itemize}
\end{frame}

\begin{frame}[fragile]{A hint on replacements in exercises and <>}
	
	If you see text enclosed in <> on our slides or in the notebook such as
	\begin{verbatim}
	<something here>
	\end{verbatim}
	it means that 'something here' should be replaced by some input from your side. (Hopefully obvious from the context)
	
	Anything else can just be copied.

\end{frame}

\section{Python introduction}
\begin{frame}[fragile]{Python introduction}
    See Jupyter notebook:\\
    \href{https://github.com/inwe-boku/lecture-scientific-computing/blob/master/lecture03-python-introduction/lecture03.ipynb}{lecture03.ipynb}
\end{frame}


\section{Homework assignment}

\begin{frame}[fragile]{Disclaimer}
    \begin{center}
        \includegraphics[height=0.6\textheight]{covid-19-expert.png}
    \end{center}

    \pause
    \textbf{Disclaimer:} Models used in this lecture are highly simplified and supposed to serve as examples
    to study Python. Don't use models as prediction. If you do your own data analysis on important
    topics, compare it with other published studies.

    {\tiny Source: \href{https://twitter.com/amrith/status/1242481942026530817}{https://twitter.com/amrith/status/1242481942026530817}}
\end{frame}

\begin{frame}[fragile]{Homework assignment}
	\begin{itemize}
		\item Fetch the latest changes from the upstream repository, to get the homework Notebook:
            {\scriptsize
            \begin{verbatim}
cd path/to/homework-scientific-computing
git pull --no-edit upstream master
git push\end{verbatim}
            }
        \item Start Jupyter and solve the exercises in the notebook:
            \href{https://github.com/inwe-boku/homework-scientific-computing/blob/master/homework03-python-introduction/homework03.ipynb}{homework03.ipynb}.
        \item Commit the notebook file and push it to your fork.
	\end{itemize}

    \bigskip
    To avoid merge conflicts, you can either commit a copy of the notebook and add your Github name
    to the filename or solve the exercises together with your group members and \verb|git pull|
    before starting to work on the notebook and \verb|git push| before the next group member starts
    working on it.
\end{frame}

\end{document}
